\documentclass[12pt]{article}

\usepackage{dsfont}
\usepackage{amsmath, amssymb,amsthm}
\usepackage{bm}
\usepackage{cite}
\usepackage[dvipsnames]{xcolor}
\usepackage{comment}
\usepackage{extarrows}
\usepackage{geometry}
\geometry{hmargin=2.5cm,vmargin=2.5cm}
\usepackage{hyperref}
\hypersetup{
	colorlinks = true,
	citecolor = cyan,
	linkcolor = magenta,
	breaklinks = true,
	linktocpage
}

\usepackage{fancyhdr}
% Define a custom style specifically for the first page
\fancypagestyle{firstpage}{
    \fancyhf{} % Clear all headers and footers
    \fancyfoot[C]{\thepage}
    \fancyfoot[R]{\today} % Put the date in the right footer
    \renewcommand{\headrulewidth}{0pt} % Remove top line
}

\usepackage{mathtools}
\DeclarePairedDelimiter{\ceil}{\lceil}{\rceil}
\usepackage{tikz}
 \usepackage{pgfplots}
%\usepackage{refcheck}


\allowdisplaybreaks

\theoremstyle{definition}
\newtheorem*{definition*}{Definition}
\newtheorem*{proposition*}{Proposition}
\newtheorem*{remark*}{Remark}
\newtheorem*{theorem*}{Theorem}
\newtheorem*{corollary*}{Corollary}
\newtheorem*{lemma*}{Lemma}
\newtheorem*{axiom*}{Axiom}
\newtheorem*{example*}{Example}
\newtheorem*{solution*}{Solution}

%\counterwithin{equation}{section}

\newcommand{\D}[2]{\prescript{}{#1}{D}_x^{#2}}
\newcommand{\I}[2]{\prescript{}{#1}{I}_x^{#2}}

%%%%%
\begin{document}
%%%%%

\begin{center}
{\bf {\Large Introduction to Fractional Calculus}}\\[1cm]
\end{center}


\begin{example*}
Compute $4!$
\end{example*}

\begin{solution*}
$4! = 1\cdot 2 \cdot 3 \cdot 4 = 24$
\end{solution*}

\begin{example*}
Using the fact that
\[
\Gamma(2) = 1! = 1,\qquad \Gamma(3) = 2! = 2,\qquad \Gamma(5/2) = \frac{3}{4}\sqrt{\pi}
\]
evaluate
\begin{itemize}
\item[(a)] $\I{0}{1/2}[x]$
\item[(b)] $\I{0}{1/2}[x^{3/2}]$
\end{itemize}
\end{example*}

\begin{solution*}
First, 
\[
\Gamma(5/2) = \Gamma(3/2+1) = \frac{3}{2}\, \Gamma(3/2) = \frac{3}{2}\cdot \frac{1}{2}\, \Gamma(1/2) = \frac{3}{4}\sqrt{\pi},
\]
\begin{itemize}
\item[(a)] Using the power rule for the Riemann--Liouville fractional derivative we get
\[
\I{0}{1/2}[x] = \frac{\Gamma(2)}{\Gamma(2+1/2)} x^{3/2} = \frac{\Gamma(2)}{\Gamma(5/2)}x^{3/2} = \frac{1}{\frac{3}{4}\sqrt{\pi}} x^{3/2} = \frac{4}{3\sqrt{\pi}}x^{3/2}.
\]
\item[(b)]Next,
\[
\I{0}{1/2}[x^{3/2}] = \frac{\Gamma(5/2)}{\Gamma(3)} x^2 = \frac{\frac{3}{4}\sqrt{\pi}}{2} x^2.
\]
\end{itemize}
Thus,
\begin{align*}
\I{0}{1/2} \I{0}{1/2}[x] &= \I{0}{1/2}\bigg[\frac{4}{3\sqrt{\pi}}x^{3/2}\bigg]
= \frac{4}{3\sqrt{\pi}} \I{0}{1/2}[x^{3/2}] \\
&= \frac{4}{3\sqrt{\pi}}\cdot \frac{\frac{3}{4}\sqrt{\pi}}{2} x^2 \\
& = \frac{x^2}{2} = \I{0}{1}[x] = \int_0^x t\, dt.
\end{align*}
\end{solution*}

\begin{example*}
Using the fact that 
\[
\Gamma\bigg(\frac{7}{2}\bigg) = \frac{15\sqrt{\pi}}{8},
\]
compute the fractional derivative of $x^2$.
\end{example*}

\begin{solution*}
Computing the half derivative of $x^2$ we obtain
\[
\D{0}{1/2}[x^2] = \frac{d}{dx}\big(\I{0}{1/2}[x^2]\big) =
 \frac{d}{dx}\bigg(\frac{\Gamma(3)}{\Gamma(7/2)}x^{5/2}\bigg) 
= \frac{5\, \Gamma(3)}{2\, \Gamma(7/2)} x^{3/2} = \frac{5}{\frac{15\sqrt{\pi}}{8}} x^{3/2} = \frac{8 x^{3/2}}{3\sqrt{\pi}}.
\]
Differentiating the result obtained
\[
\D{0}{1/2}\bigg[\frac{8 x^{3/2}}{3\sqrt{\pi}}\bigg] = 
\frac{8}{3\sqrt{\pi}} \frac{d}{dx}\big(\I{0}{1/2}[x^{3/2}]\big)
=\frac{8}{3\sqrt{\pi}} \cdot \frac{\frac{3}{4}\sqrt{\pi}}{2} \frac{d}{dx}(x^2) = 2 x.
\]
As expected,
\[
\D{0}{1/2}\circ \D{0}{1/2}[x^2] = \D{0}{1}[x^2] = \frac{d(x^2)}{dx}.
\]
\end{solution*}

\end{document}